\section{Conclusion}
\label{sec:conclusion}

\subsection{Contribution}

This study made three principal contributions to the field of text-based emotion classification:

\begin{enumerate}[nosep]
    \item \textbf{Unified benchmarking framework.} A controlled comparison of SVM, XGBoost, CNN, and BiLSTM was conducted on a 416,809-sample Twitter corpus under identical preprocessing, stratified splits, and macro-F1 evaluation --- addressing the methodological inconsistencies that have hindered direct cross-study comparisons in prior work.
    \item \textbf{Error taxonomy.} Systematic confusion matrix analysis identified Joy$\leftrightarrow$Love (valence conflation) and Fear$\rightarrow$Surprise (arousal ambiguity) as architecture-independent misclassification patterns, establishing that semantic overlap --- rather than model capacity --- defines the current classification boundary.
    \item \textbf{End-to-end deployment pipeline.} The best-performing classical model (SVM) was integrated into a production-ready Flask REST API with Docker containerisation and CI/CD automation, demonstrating that high-accuracy emotion prediction can be delivered at sub-millisecond inference latency.
\end{enumerate}

The complete codebase --- preprocessing scripts, model training notebooks, and deployment configuration --- is publicly available to support reproducibility and future extension.

\subsection{Summary}

This study set out to answer the question: \emph{How do traditional ML models (SVM, XGBoost) compare with deep learning models (CNN, BiLSTM) for multi-class emotion classification under unified experimental conditions?}

The results demonstrate that BiLSTM achieved the highest test accuracy (93.93\%, macro-F1 0.9391), followed by CNN (92.54\%), SVM (91.80\%), and XGBoost (90.79\%). The consistent ranking BiLSTM $>$ CNN $>$ SVM $>$ XGBoost confirms that bidirectional contextual encoding provides a measurable advantage over local feature extraction and bag-of-words representations for emotion classification. Confusion matrix analysis revealed that Joy$\leftrightarrow$Love and Fear$\rightarrow$Surprise constituted the dominant misclassification pairs across all four architectures, confirming that semantic overlap is the primary classification bottleneck.

From a deployment perspective, BiLSTM's ${\sim}23\times$ higher inference latency relative to SVM made it impractical for latency-sensitive production use. SVM was therefore selected and deployed as a Flask-based REST API containerised with Docker, demonstrating that a classical model with TF-IDF features can deliver reliable real-time predictions at minimal computational cost.

\subsection{Limitations}

Several limitations should be acknowledged. The dataset is drawn exclusively from Twitter, and the informal, character-limited style of tweets may not generalise to longer-form text domains such as product reviews or clinical notes. Crowd-sourced emotion labels introduce annotation noise, particularly for ambiguous or sarcastic posts. The six-class single-label taxonomy cannot capture mixed emotions (\eg, bittersweet posts combining Joy and Sadness). Finally, transformer-based models (BERT, RoBERTa) were excluded due to computational constraints, leaving the comparison incomplete with respect to the current state of the art.

\subsection{Future Work}

Four directions are identified for future investigation:
\begin{enumerate}[nosep]
    \item Fine-tuning transformer-based models (BERT, RoBERTa) on the same corpus to benchmark against the baselines established in this study.
    \item Cross-domain evaluation on non-Twitter corpora (Reddit, product reviews) to assess the generalisability of the trained models.
    \item Development of a real-time emotion monitoring dashboard to extend the deployed API for streaming social media analysis.
    \item Explainability analysis using SHAP and LIME to identify token-level features that discriminate between semantically similar emotion pairs.
\end{enumerate}
